\documentclass[11pt]{article}

\usepackage[margin=1in]{geometry}
\usepackage{amsmath, amssymb, amsthm}
\usepackage{mathtools}
\usepackage{bm}
\usepackage{hyperref}
\usepackage{xcolor}
\usepackage{enumitem}

\newtheorem{theorem}{Theorem}
\newtheorem{lemma}[theorem]{Lemma}
\newtheorem{proposition}[theorem]{Proposition}
\newtheorem{corollary}[theorem]{Corollary}
\theoremstyle{definition}
\newtheorem{definition}[theorem]{Definition}
\newtheorem{assumption}[theorem]{Assumption}
\theoremstyle{remark}
\newtheorem{remark}[theorem]{Remark}

\newcommand{\R}{\mathbb{R}}
\newcommand{\E}{\mathbb{E}}
\newcommand{\Var}{\mathrm{Var}}
\newcommand{\Cov}{\mathrm{Cov}}
\newcommand{\KL}{\mathrm{KL}}
\newcommand{\TV}{\mathrm{TV}}
\newcommand{\D}{\mathcal{D}}
\newcommand{\Loss}{\mathcal{L}}
\newcommand{\eps}{\varepsilon}
\newcommand{\old}{\mathrm{old}}
\DeclareMathOperator*{\argmax}{arg\,max}

\title{Variational-Dropout Stochasticity for GRPO Fine-Tuning of \\
       Latent-Reasoning Language Models \\[4pt]
       \large (revised)}
\author{Wooil Jung}
\date{\today}

\begin{document}
\maketitle

\begin{abstract}
Latent-reasoning language models such as \textsc{Coconut} (Hao et al., 2024)
replace discrete chain-of-thought tokens with continuous hidden states fed
recurrently into the next reasoning step.  Because the latent recurrence is
deterministic given the input, na\"ive policy-gradient methods such as GRPO
cannot generate diverse rollouts in latent space.  We propose to use a small
dropout rate $p \approx 0.1$ as the source of stochasticity during rollout, and
to \emph{replay the same dropout mask} during the policy update so that the
latent trajectory observed at rollout time is reproduced exactly.  Following
Gal \& Ghahramani (2016), we further \emph{share the mask across all latent
steps within one rollout}, so each rollout corresponds to a single posterior
sample $\tilde\theta(\xi)$ from the variational dropout distribution.  We show
that, under this construction, dropout noise is exogenous, the per-mask
advantage $A^{(k)} = r^{(k)} - \mu_r$ (without standard-deviation
normalization, following Liu et al.\ 2024) yields an unbiased policy-gradient
estimator on the dropout-marginalized policy, and replaying the mask provides
common-random-numbers variance reduction over freshly resampling at update
time.  We also give an inductive argument that the surrogate gradient is
well defined at every latent depth $T$.  A reference implementation is
provided in \texttt{coconut\_variational.py} and
\texttt{run\_grpo\_dropout\_variational.py}.
\end{abstract}

% =====================================================================
\section{Revision summary}
% =====================================================================

This revision corrects two errors and one ambiguity in the previous draft,
and documents four implementation refinements adopted during empirical
debugging.

\begin{itemize}[topsep=2pt,itemsep=2pt]
\item \textbf{Hole 1 (advantage normalization).}  The original draft used
      the GRPO advantage
      $A^{(k)} = (r^{(k)} - \mu_r)/(\sigma_r + \delta)$
      and claimed unbiasedness in Theorem 4.4.  The proof handled only the
      mean baseline $\mu_r$ and silently ignored $\sigma_r$.  Because
      $\sigma_r$ is a nonlinear function of every $r^{(j)}$ in the group
      (including $r^{(k)}$ itself), the control-variate identity does not
      extend to dividing by it, and the estimator is biased.  This is the
      same observation made independently by Liu et al.\ (2024)
      (``Dr.\ GRPO'').  In this revision we drop the $\sigma_r$
      normalization, use $A^{(k)} = r^{(k)} - \mu_r$, and recover
      $(1 - 1/K)\,\nabla J(\theta)$ exactly.
\item \textbf{Hole 2 (CRN baseline).}  The previous Proposition 4.5
      compared mask replay to an estimator that re-used the rollout's
      reward $A(\xi,y)$ but evaluated the score function on a freshly
      sampled mask $\xi'$.  Because $y$ was not generated under $\xi'$,
      the resampling estimator is malformed: it is biased and the mean
      equality assumed in the proof fails.  This revision replaces the
      baseline with a \emph{fresh-rollout} estimator
      $\hat g_{\mathrm{indep}} = A(\xi',y')\,\nabla \log \pi_\theta(y'\mid x,\xi')$
      where $(\xi',y')$ is drawn afresh at update time.  The CRN argument
      then goes through.
\item \textbf{MC-dropout (formerly ambiguous).}  The previous draft drew
      independent per-step masks $m_t$, which is incompatible with the
      Gal--Ghahramani interpretation in which a single posterior sample
      $\theta'$ is reused throughout a forward pass.  The implementation
      in \texttt{coconut\_variational.py} now shares one mask
      $m_t \equiv m$ across all latent recurrence steps via an
      RNG-restore mechanism.  Section~\ref{sec:mc-dropout} of this
      revision relies on that shared-mask construction and gives the
      Bayesian model-average reading rigorously.
\item \textbf{Implementation refinements (V4--V7).}  Section~\ref{sec:impl}
      documents four changes adopted to stabilize training in practice:
      (V4) a Huberized $k_3$ KL estimator that keeps tail gradients bounded
      but nonzero, replacing the standard log-ratio clamp;
      (V5) an annealed KL coefficient $\beta(t) = \beta_0 \cdot \eta_t/\eta_{\max}$
      that shrinks the trust-region pull alongside the learning-rate
      schedule;
      (V6) sequence-mean loss aggregation in place of token-mean, matching
      Algorithm~\ref{alg:dropout-grpo} and removing length-weighting
      artifacts;
      (V7) a group-level accuracy filter
      $\alpha_{\min} \le \mu_{r}^{(g)} \le \alpha_{\max}$ that drops
      uninformative groups before any replay or policy update, with
      collective DDP coordination so that empty global batches skip the
      step entirely (no optimizer step, no scheduler step).
\end{itemize}

The well-definedness lemma, the smoothness assumption, and the original
practical discussion are unchanged in substance.

% =====================================================================
\section{Introduction and motivation}
% =====================================================================

Group Relative Policy Optimization (GRPO; Shao et al., 2024) has become a
standard reinforcement-learning objective for fine-tuning reasoning LLMs on
math benchmarks.  Its appeal is that the per-prompt advantage is computed
purely from a group of $K$ rollouts and a verifier reward, with no learned
critic.  GRPO requires that the $K$ rollouts \emph{differ}, because the
group-mean baseline,
\begin{equation}
A^{(k)} \;=\; r^{(k)} - \mu_r,
\qquad
\mu_r \;=\; \tfrac1K\sum_{k=1}^{K} r^{(k)},
\label{eq:grpo-advantage}
\end{equation}
is identically zero whenever all rollouts agree.  In a standard token-level
LLM, diversity is supplied automatically by sampling from the next-token
distribution.

\textsc{Coconut}-style latent reasoning, by contrast, does not sample.  Hidden
states $h_t \in \R^d$ are fed back into the model as the embedding of the next
``token'' for $t = 1, \dots, T$ until a stopping criterion fires; only after
the latent phase ends does the model emit discrete answer tokens
$y_1, \dots, y_L$.  Conditional on the prompt $x$, the latent phase is a
\emph{deterministic} function of the parameters $\theta$.  This breaks GRPO:
all $K$ rollouts share the same $h_T$ and (for greedy decoding) the same $y$.

We propose to inject stochasticity by enabling dropout during rollout, with
one dropout mask drawn per rollout and held constant across all latent
recurrence steps within that rollout (variational dropout, Gal \& Ghahramani
2016).  This yields per-rollout latent trajectories that depend on a
Bernoulli mask $\xi \sim p(\xi)$ which is independent of $\theta$, so GRPO
can grade rollouts relative to one another.  At update time we re-mount the
\emph{same} mask so that the gradient is computed against the trajectory
that was actually graded.

\paragraph{Note on advantage normalization.}
We deliberately omit the standard-deviation normalization
$1/(\sigma_r + \delta)$ that appears in the original GRPO formulation.  As
argued by Liu et al.\ (2024) and discussed in Section~\ref{sec:revision},
dividing by a within-group statistic of the rewards introduces bias because
$\sigma_r$ depends on every individual $r^{(k)}$.  Empirically, the
normalization can be useful as an adaptive learning-rate trick, but it is
not required for correctness and is omitted here for theoretical clarity.

% =====================================================================
\section{Preliminaries}
% =====================================================================

\subsection{Coconut recurrence with variational dropout}
\label{sec:coconut-recurrence}

Let $f_\theta : \R^d \to \R^d$ denote one Transformer step of the
latent-reasoning model with parameters $\theta \in \R^D$.  Standard
\textsc{Coconut} unrolls
\begin{equation}
h_t \;=\; f_\theta(h_{t-1}),
\qquad
h_0 \;=\; \mathrm{embed}_\theta(x),
\qquad t = 1,\dots,T.
\end{equation}
With dropout at rate $p$ and keep-probability $q := 1-p$, draw a
\emph{single} Bernoulli mask $m \in \{0,1\}^D$ with
$m \sim \mathrm{Bern}(q)^{\otimes D}$ per rollout, and define rescaled
parameters
\begin{equation}
\tilde\theta(\xi) \;:=\; \theta \odot m / q,
\qquad
\E_\xi\!\big[m/q\big] = \mathbf 1.
\label{eq:dropout-rescale}
\end{equation}
The dropout-perturbed recurrence reuses the \emph{same} $\tilde\theta(\xi)$
at every step:
\begin{equation}
h_t \;=\; f_{\tilde\theta(\xi)}(h_{t-1}),
\qquad
\xi \sim p(\xi),
\qquad t = 1,\dots,T.
\label{eq:dropout-recurrence}
\end{equation}
This is variational dropout in the sense of Gal \& Ghahramani (2016): one
posterior sample $\tilde\theta(\xi)$ is reused across the entire recurrence,
not redrawn per step.  Once the latent phase finishes, answer tokens are
sampled autoregressively,
\begin{equation}
y_\ell \;\sim\; \pi_\theta(\cdot \mid x, h_T, y_{<\ell}),
\qquad \ell = 1, \dots, L,
\end{equation}
optionally with dropout enabled here as well; we absorb any such sampling
randomness into $\xi$ by reparameterizing the categorical samples with
Gumbel noise when convenient.  A scalar reward $r(y, x) \in [0,1]$ is provided
by an external verifier (e.g.\ exact-match on \textsc{GSM8K}).

\begin{remark}[Implementation]
The shared mask is realized in \texttt{coconut\_variational.py} by
capturing the CUDA + CPU RNG state at the start of each chain and
restoring it at the start of each latent recurrence step
(\texttt{\_forward\_base} branch on \texttt{inputs\_embeds.shape[1] == 1}).
Constant-shape dropout layers (residual / MLP / embedding) therefore draw
bit-identical Bernoullis at every step.  Variable-shape attention dropout,
whose mask shape grows with kv-cache length, is checked by an auditor
that compares the prefix overlap; in PyTorch's Philox-based dropout
kernel the prefix is determined by the seed alone, so prefix masks at
steps $t$ and $t-1$ remain bit-identical on the overlap.  See the file
header for caveats on architectures (Qwen2 / LLaMA / Mistral / Gemma)
that ship without residual / MLP dropout.
\end{remark}

\subsection{Augmented policy}
The induced distribution over $(\xi, y)$ given $x$ is
\begin{equation}
\Pi_\theta(\xi, y \mid x)
   \;=\; p(\xi) \, \pi_\theta(y \mid x, \xi),
\qquad
\pi_\theta(y \mid x, \xi)
  := \prod_{\ell=1}^{L} \pi_\theta\!\big(y_\ell \,\big|\, x, h_T(\theta,\xi,x), y_{<\ell}\big).
\label{eq:augmented-policy}
\end{equation}
We treat $\Pi_\theta$ as the policy that GRPO optimizes; the dropout mask
is part of the action.  The marginal policy
\begin{equation}
\bar\pi_\theta(y \mid x)
  \;:=\; \E_{\xi \sim p}\!\big[\pi_\theta(y \mid x, \xi)\big]
\label{eq:marginal-policy}
\end{equation}
is the object whose expected reward we ultimately wish to maximize.

\subsection{GRPO objective}

For a group of $K$ rollouts $\{(\xi^{(k)}, y^{(k)})\}_{k=1}^K$ generated
under parameters $\theta_\old$, with rewards $r^{(k)} = r(y^{(k)}, x)$,
define the advantage $A^{(k)}$ via~\eqref{eq:grpo-advantage} and the
per-token importance ratio
\begin{equation}
\rho_\theta^{(k)}(\xi, y)
  \;:=\;
\frac{\pi_\theta(y \mid x, \xi)}{\pi_{\theta_\old}(y \mid x, \xi)}.
\end{equation}
The clipped GRPO surrogate is
\begin{equation}
\Loss^{\mathrm{GRPO}}(\theta)
  \;=\;
\E_{x \sim \D}\,
\frac{1}{K}\sum_{k=1}^{K}
  \min\!\Big(
    \rho_\theta^{(k)}\, A^{(k)}, \;
    \mathrm{clip}(\rho_\theta^{(k)},\,1-\eps,\,1+\eps)\,A^{(k)}
  \Big)
  \;-\;\beta\,\KL\!\big(\pi_{\theta}\,\big\|\,\pi_{\mathrm{ref}}\big).
\label{eq:grpo-loss}
\end{equation}

% =====================================================================
\section{Method: dropout as exogenous noise with mask replay}
\label{sec:method}
% =====================================================================

The proposed procedure is:
\begin{enumerate}[topsep=2pt,itemsep=2pt]
\item For each prompt $x$, draw $K$ independent dropout masks
      $\xi^{(1)}, \dots, \xi^{(K)} \sim p(\xi)$.  Each mask is a single
      sample reused across all $T$ latent steps within its rollout.
\item For each $k$, run the latent-reasoning recurrence
      \eqref{eq:dropout-recurrence} with mask $\xi^{(k)}$ (RNG seed stored
      on device), decode answer tokens $y^{(k)}$, compute reward
      $r^{(k)}$ and advantage $A^{(k)}$ via~\eqref{eq:grpo-advantage}.
\item During the policy update, recompute
      $\pi_\theta(y^{(k)}\mid x,\xi^{(k)})$ with the \emph{same} stored
      mask $\xi^{(k)}$, so that $h_T$ at update time matches $h_T$ at
      rollout time exactly when $\theta = \theta_\old$.  Form the GRPO
      loss~\eqref{eq:grpo-loss} and step $\theta$.
\end{enumerate}

The construction rests on three observations.

\paragraph{(i) Independence of $\xi$ from $\theta$.}
The dropout mask $\xi$ is sampled from a fixed distribution $p(\xi)$ that
does not depend on $\theta$.  Hence, in~\eqref{eq:augmented-policy},
$\nabla_\theta \log p(\xi) = 0$, and
\begin{equation}
\nabla_\theta \log \Pi_\theta(\xi, y \mid x)
  \;=\;
\nabla_\theta \log \pi_\theta(y \mid x, \xi).
\label{eq:score-reduces}
\end{equation}
The gradient signal therefore comes entirely from the $\xi$-conditional log
likelihood of the answer tokens, which is differentiable through the
unrolled latent path.

\paragraph{(ii) Stochastic-computation-graph view.}
With $\xi$ fixed, $h_T = h_T(\theta, \xi, x)$ is a deterministic, almost
everywhere differentiable function of $\theta$.  The estimator
$A^{(k)}\,\nabla_\theta \log \pi_\theta(y^{(k)}\mid x,\xi^{(k)})$ is a
score-function estimator on $\Pi_\theta$ with the dropout draw treated as
an exogenous input (Schulman et al.\ 2015, ``Gradient estimation using
stochastic computation graphs'').  Because $y$ is categorical we cannot
apply the pathwise / reparameterization trick at the answer-token boundary;
the gradient through the latent dynamics nonetheless flows by ordinary
backpropagation through the unrolled recurrence.

\paragraph{(iii) Mask replay = common random numbers.}
At update time we evaluate $\pi_\theta$ on the same $\xi^{(k)}$ that was
used at rollout time.  This couples the rollout-time and update-time
trajectories so that the only source of difference between them is the
parameter update $\theta - \theta_\old$ itself.  At $\theta = \theta_\old$
the importance ratio $\rho_{\theta_\old}^{(k)}$ is identically $1$ and
$h_T$ at update time exactly reproduces $h_T$ at rollout time --- a
discrepancy that would otherwise contribute additional variance to the
estimator.  Section~\ref{sec:crn} formalizes this CRN claim against a
fresh-rollout baseline.

% =====================================================================
\section{Theoretical analysis}
\label{sec:revision}
% =====================================================================

\subsection{Well-definedness of the surrogate gradient (induction on $T$)}

\begin{assumption}\label{ass:smooth}
For every fixed mask $\xi$ and almost every $\theta \in \R^D$ in Lebesgue
measure, the map $\theta \mapsto f_{\tilde\theta(\xi)}(h)$ is differentiable
in $\theta$ for almost every $h$, and the answer-token log-likelihood
$\log \pi_\theta(y \mid x, h_T, y_{<\ell})$ is differentiable in both
$\theta$ and $h_T$.
\end{assumption}

This is satisfied by every standard Transformer block: the dropout
rescaling $\theta \odot m/q$ is linear in $\theta$, and the remaining
nonlinearities (LayerNorm, GELU, softmax) are almost everywhere
differentiable.

\begin{lemma}[Well-defined latent gradient at depth $T$]
\label{lem:induction}
Under Assumption~\ref{ass:smooth}, for every $T \ge 0$ and every $\xi$,
the Jacobian $\partial h_T / \partial \theta$ exists almost everywhere in
$\theta$ and is given by the recursion
\begin{equation}
\frac{\partial h_t}{\partial \theta}
  \;=\;
\frac{\partial f_{\tilde\theta}}{\partial \theta}\bigg|_{h_{t-1}}
  \cdot
\frac{\partial \tilde\theta}{\partial \theta}
  \;+\;
\frac{\partial f_{\tilde\theta}}{\partial h}\bigg|_{h_{t-1}}
  \cdot
\frac{\partial h_{t-1}}{\partial \theta},
\qquad
\frac{\partial h_0}{\partial \theta}
  \;=\;
\frac{\partial \,\mathrm{embed}_\theta(x)}{\partial \theta},
\label{eq:latent-jacobian}
\end{equation}
where $\partial \tilde\theta / \partial \theta = \mathrm{diag}(m)/q$ is the
dropout-rescaled identity.  Consequently
$\nabla_\theta \log \pi_\theta(y \mid x, \xi)$ is well defined.
\end{lemma}

\begin{proof}
By induction on $t$.

\emph{Base case} ($t=0$).  $h_0 = \mathrm{embed}_\theta(x)$ is
differentiable in $\theta$ by assumption, giving the stated initial
Jacobian.

\emph{Inductive step.}  Assume $\partial h_{t-1}/\partial \theta$ exists.
Write $h_t = f(\tilde\theta, h_{t-1})$, viewing $f$ as a function of two
arguments.  By the multivariate chain rule and Assumption~\ref{ass:smooth},
\[
\frac{\partial h_t}{\partial \theta}
  \;=\;
\frac{\partial f}{\partial \tilde\theta}
  \cdot
\frac{\partial \tilde\theta}{\partial \theta}
  \;+\;
\frac{\partial f}{\partial h_{t-1}}
  \cdot
\frac{\partial h_{t-1}}{\partial \theta},
\]
and $\partial \tilde\theta/\partial \theta = \mathrm{diag}(m)/q$ from
\eqref{eq:dropout-rescale}.  Both terms exist (the second by the
inductive hypothesis), giving~\eqref{eq:latent-jacobian}.

Finally, since $\log \pi_\theta(y \mid x, h_T, y_{<\ell})$ is
differentiable in both $\theta$ and $h_T$, one more chain-rule
application yields
$\nabla_\theta \log \pi_\theta(y \mid x, \xi)
   = \sum_\ell [\partial_\theta \log \pi_\theta]_{\text{direct}}
              + [\partial_{h_T}\log \pi_\theta]\cdot \partial h_T/\partial\theta$.
\end{proof}

\begin{remark}
The recursion~\eqref{eq:latent-jacobian} is what \texttt{loss.backward()}
implements when the dropout mask is held constant via RNG restore: PyTorch's
\texttt{torch.nn.functional.dropout} treats the rescaled mask as a constant
during backprop, which is exactly $\partial \tilde\theta/\partial \theta =
\mathrm{diag}(m)/q$.
\end{remark}

\subsection{Unbiasedness of the dropout-GRPO surrogate}
\label{sec:unbiased}

Throughout this subsection we condition on a single prompt $x$ and treat
the $K$ rollouts as i.i.d.\ draws from
$\Pi_{\theta_\old}(\cdot \mid x) = p(\xi)\,\pi_{\theta_\old}(y\mid x,\xi)$.

\begin{theorem}[Unbiased policy-surrogate gradient]
\label{thm:unbiased}
Let $J(\theta) := \E_{x\sim\D}\,\E_{(\xi,y)\sim\Pi_\theta(\cdot\mid x)}[r(y,x)]$
be the expected reward under the augmented policy~\eqref{eq:augmented-policy}.
Define the surrogate part of the GRPO loss (excluding the KL term) as
$\Loss^{\mathrm{surr}}(\theta) := \Loss^{\mathrm{GRPO}}(\theta)
  + \beta\,\KL(\pi_\theta\,\|\,\pi_{\mathrm{ref}})$.
Under the advantage definition $A^{(k)} = r^{(k)} - \mu_r$
in~\eqref{eq:grpo-advantage},
\begin{equation}
\nabla_\theta \Loss^{\mathrm{surr}}(\theta)\Big|_{\theta=\theta_\old}
  \;=\;
\Big(1 - \tfrac{1}{K}\Big)\,\nabla_\theta J(\theta)\Big|_{\theta_\old}.
\label{eq:unbiased}
\end{equation}
Equivalently,
$\nabla_\theta \Loss^{\mathrm{GRPO}}(\theta)\big|_{\theta_\old}
   = (1 - 1/K)\,\nabla J(\theta_\old)
   - \beta\,\nabla_\theta \KL(\pi_\theta\,\|\,\pi_{\mathrm{ref}})\big|_{\theta_\old}$.
\end{theorem}

\begin{proof}
At $\theta = \theta_\old$ we have $\rho_{\theta_\old}^{(k)} \equiv 1$, so
the clip is inactive and the gradient of the $\min$ reduces to the
gradient of the unclipped ratio:
\[
\nabla_\theta \rho_\theta^{(k)}\big|_{\theta_\old}
   \;=\;
\nabla_\theta \log \pi_\theta(y^{(k)}\mid x,\xi^{(k)})\big|_{\theta_\old}.
\]
Therefore, with $s^{(k)} := \nabla_\theta \log \pi_\theta(y^{(k)}\mid x,\xi^{(k)})$
evaluated at $\theta_\old$,
\begin{equation}
\nabla_\theta \Loss^{\mathrm{surr}}\big|_{\theta_\old}
  \;=\;
\E_{x,\,\{(\xi^{(k)},y^{(k)})\}}\!\left[\frac{1}{K}\sum_{k=1}^K
   (r^{(k)} - \mu_r)\,s^{(k)}\right].
\label{eq:proof-step1}
\end{equation}
We expand the baseline term.  Because rollouts are i.i.d.\ given $x$,
\[
\E[\mu_r\,s^{(k)}]
  \;=\;
\frac{1}{K}\sum_{j=1}^K \E[r^{(j)}\,s^{(k)}]
  \;=\;
\frac{1}{K}\E[r^{(k)}\,s^{(k)}]
  \;+\;
\frac{1}{K}\sum_{j\neq k}\E[r^{(j)}]\,\E[s^{(k)}].
\]
The score function has zero mean: $\E[s^{(k)}] = \E_\xi \E_{y\mid \xi}
[\nabla\log\pi_\theta(y\mid x,\xi)] = \nabla \E_y[1] = 0$.  Hence
\[
\E[\mu_r\,s^{(k)}] \;=\; \tfrac{1}{K}\E[r^{(k)}\,s^{(k)}],
\]
and substituting back,
\begin{align*}
\E[(r^{(k)} - \mu_r)\,s^{(k)}]
  &= \E[r^{(k)}\,s^{(k)}] - \tfrac{1}{K}\E[r^{(k)}\,s^{(k)}] \\
  &= \big(1 - \tfrac{1}{K}\big)\,\E[r^{(k)}\,s^{(k)}].
\end{align*}
By~\eqref{eq:score-reduces},
$s^{(k)} = \nabla_\theta\log\Pi_\theta(\xi^{(k)},y^{(k)}\mid x)$, so
$\E[r^{(k)}\,s^{(k)}] = \nabla_\theta J(\theta)|_{\theta_\old}$ is the
standard REINFORCE policy gradient on the augmented policy.  Averaging
over $k$ preserves this and yields~\eqref{eq:unbiased}.

The KL term contributes
$-\beta\,\nabla_\theta\KL(\pi_\theta\|\pi_{\mathrm{ref}})|_{\theta_\old}$,
which is in general nonzero (it vanishes only if
$\pi_{\mathrm{ref}} = \pi_{\theta_\old}$, e.g.\ if the reference is reset
each step).  This term is the regularizer's contribution and has known
gradient.
\end{proof}

\begin{corollary}[Optimization on the marginal policy]
\label{cor:marginal}
The estimator targets $\nabla J(\theta)$, the expected reward of the
mask-marginalized policy~\eqref{eq:marginal-policy}, since
$\E_\xi[\pi_\theta(y\mid x,\xi)] = \bar\pi_\theta(y\mid x)$ implies
$\E_{\xi,y\sim\Pi_\theta}[r] = \E_{y\sim\bar\pi_\theta}[r]$.  In other
words, optimizing $\Pi_\theta$ is exactly optimizing expected reward under
the dropout ensemble.
\end{corollary}

\begin{remark}[Why $\sigma_r$ is dropped]
\label{rem:sigma}
If we instead used $A^{(k)} = (r^{(k)} - \mu_r)/(\sigma_r + \delta)$, the
calculation in the proof above would require expanding
$\E[(r^{(k)} - \mu_r)/(\sigma_r+\delta)\cdot s^{(k)}]$.  Because
$\sigma_r$ is a nonlinear function of every $r^{(j)}$ (including
$r^{(k)}$), the factorization
$\E[r^{(j)}\,s^{(k)}] = \E[r^{(j)}]\E[s^{(k)}]$ used for $j \neq k$
breaks: $\sigma_r$ couples $s^{(k)}$ to all rewards through a non-affine
function.  The resulting estimator overweights low-variance prompts (where
$1/\sigma_r$ is large) and underweights high-variance ones, distorting the
per-prompt gradient direction.  Liu et al.\ (2024) report this bias
empirically and trace prompt-difficulty drift to it.  We therefore drop
the $\sigma_r$ normalization throughout, at the cost of losing automatic
adaptive scaling --- a tradeoff that recovers theoretical correctness in
exchange for one extra learning-rate hyperparameter.
\end{remark}

\subsection{Variance reduction via mask replay (CRN)}
\label{sec:crn}

We compare two unbiased estimators of $\nabla_\theta J$ at parameter
$\theta = \theta_\old + \Delta$ for small $\Delta$.

\begin{itemize}[topsep=2pt,itemsep=2pt]
\item \textbf{Mask-replay (CRN) estimator.}  The same $\xi$ used at
      rollout time is reused at update time.  Conditional on $x$ and on
      the rollout draw $(\xi, y) \sim \Pi_{\theta_\old}(\cdot\mid x)$,
      \begin{equation}
      \hat g_{\mathrm{CRN}}(\theta;\xi,y)
        \;:=\;
      A(\xi,y)\,\nabla_\theta \log \pi_\theta(y\mid x,\xi).
      \end{equation}
\item \textbf{Fresh-rollout estimator.}  At update time we draw a new
      mask $\xi'$ and sample $y' \sim \pi_{\theta_\old}(\cdot\mid x,\xi')$
      from a fresh rollout, recompute its advantage
      $A(\xi',y')$ against an independent group, and form
      \begin{equation}
      \hat g_{\mathrm{fresh}}(\theta;\xi',y')
        \;:=\;
      A(\xi',y')\,\nabla_\theta \log \pi_\theta(y'\mid x,\xi').
      \end{equation}
\end{itemize}

\begin{remark}[On the previous draft's baseline]
The original Proposition 4.5 used
$A(\xi,y)\,\nabla_\theta\log\pi_\theta(y\mid x,\xi')$ with $\xi'\perp\xi$
fresh and $y$ from the rollout.  This estimator is biased --- $y$ was not
sampled under $\xi'$, so the score-function identity does not hold ---
and a CRN argument relative to it is meaningless.  The fresh-rollout
baseline above is the correct counterfactual: ``what if you had not
replayed the mask, but instead drawn a brand-new rollout at update
time.''
\end{remark}

\begin{proposition}[CRN variance reduction]
\label{prop:crn}
Both $\hat g_{\mathrm{CRN}}$ and $\hat g_{\mathrm{fresh}}$ are unbiased
estimators of $\nabla J(\theta)$ at $\theta = \theta_\old$ (modulo the
$(1-1/K)$ scale of Theorem~\ref{thm:unbiased}).  Suppose the function
$\theta \mapsto A(\xi,y)\,\nabla_\theta\log\pi_\theta(y\mid x,\xi)$ is
$L$-Lipschitz in $\theta$ for $\Pi_{\theta_\old}$-almost every $(\xi,y)$,
and that the difference
\[
\Delta(\theta;\xi,y) \;:=\;
   A(\xi,y)\,\nabla_\theta\log\pi_\theta(y\mid x,\xi)
   - A(\xi,y)\,\nabla_\theta\log\pi_{\theta_\old}(y\mid x,\xi)
\]
satisfies $\|\Delta(\theta;\xi,y)\| \le L\|\theta - \theta_\old\|$.  Then
\begin{equation}
\Var[\hat g_{\mathrm{CRN}}(\theta)]
  \;\le\;
\Var[\hat g_{\mathrm{fresh}}(\theta)]
  \;+\;
2L\,\|\theta - \theta_\old\|\,\sqrt{\Var[\hat g_{\mathrm{CRN}}(\theta_\old)]}
  \;+\;
L^2\|\theta - \theta_\old\|^2.
\end{equation}
In particular, in the trust-region regime $\|\theta - \theta_\old\| \to 0$
the CRN variance lower-bounds the fresh-rollout variance up to
$O(\|\theta - \theta_\old\|)$ slack.  At $\theta = \theta_\old$ exactly,
$\hat g_{\mathrm{CRN}}(\theta_\old)$ has the same per-sample value
distribution as $\hat g_{\mathrm{fresh}}(\theta_\old)$, but reuses
$(\xi,y)$ across calls, eliminating the across-call variance contribution
from resampling.
\end{proposition}

\begin{proof}[Proof sketch]
Write each estimator as a value at $\theta_\old$ plus a perturbation:
$\hat g_{\mathrm{CRN}}(\theta) = \hat g_{\mathrm{CRN}}(\theta_\old) +
   \Delta(\theta;\xi,y)$ and analogously for $\hat g_{\mathrm{fresh}}$
with $(\xi',y')$.  The two zero-order terms have the same distribution
because $(\xi,y)$ and $(\xi',y')$ are both drawn from
$\Pi_{\theta_\old}(\cdot\mid x)$.  The variance bound follows from
Cauchy--Schwarz on the cross term and the Lipschitz hypothesis.  This is
the standard CRN argument (Glasserman \& Yao, 1992, Thm 3.1) specialized
to a score-function estimator coupled across $\theta$ values via shared
exogenous noise.
\end{proof}

\begin{remark}[Practical reading]
The bound says: at $\theta_\old$ both estimators have identical variance,
and CRN's advantage shows up cumulatively across mini-epochs as $\theta$
drifts from $\theta_\old$.  The fresh-rollout baseline pays an additional
across-mask variance term that CRN avoids by construction.  In a
trust-region regime where $\|\theta - \theta_\old\|$ is kept small by the
PPO clip, this advantage is the dominant variance reduction.
\end{remark}

\subsection{MC-dropout / Bayesian model-average reading}
\label{sec:mc-dropout}

Following Gal \& Ghahramani (2016), the shared-mask construction
\eqref{eq:dropout-recurrence} is interpretable as a single posterior
sample from a structured variational distribution $q_\theta(\theta')$ over
parameters: $\tilde\theta(\xi) = \theta\odot m(\xi)/q$ with
$m \sim \mathrm{Bern}(q)^{\otimes D}$.  Crucially, because the mask is held
constant across all $T$ latent steps, the entire latent recurrence is
evaluated at a single $\tilde\theta(\xi)$, matching the variational
inference derivation.  The marginal policy~\eqref{eq:marginal-policy} is
then exactly a Bayesian model-average policy:
\begin{equation}
\bar\pi_\theta(y\mid x)
  \;=\;
\E_{\theta'\sim q_\theta}[\pi_{\theta'}(y\mid x)].
\label{eq:bma}
\end{equation}
Theorem~\ref{thm:unbiased} and Corollary~\ref{cor:marginal} together state
that dropout-GRPO performs ensemble-aware policy optimization: each
gradient step is the Monte Carlo average of per-ensemble-member policy
gradients, with the group-mean baseline acting as a control variate.

\begin{remark}[Implementation caveats]
The Bayesian reading holds cleanly for layers whose dropout mask shape
does not depend on the step index $t$ (residual / MLP / embedding
dropout), because RNG restore reproduces the mask bit-identically.  For
attention dropout in modern decoders (Qwen2 / LLaMA / Mistral / Gemma)
the kv-cache length grows with $t$, so the mask shape varies; on the
shape-overlap prefix the mask is still bit-identical (PyTorch Philox
kernel property), but the suffix differs.  This is a per-layer
approximation of variational dropout; the proof framework in
Sections~\ref{sec:unbiased}--\ref{sec:crn} is unaffected because it only
requires $\xi$ to be exogenous and replayable.  The
\texttt{\_DropoutMaskAuditor} class in \texttt{coconut\_variational.py}
verifies the bit-identity property at runtime.
\end{remark}

% =====================================================================
\section{Algorithm}
% =====================================================================

\begin{figure}[h]
\centering
\fbox{\begin{minipage}{0.92\linewidth}
\textbf{Algorithm 1: Dropout-GRPO for latent-reasoning LLMs (revised).}\\[2pt]
\textit{Inputs:} base policy $\pi_{\theta_\old}$, prompts $\D$, group size
$K$, dropout rate $p$, latent depth $T$, clip $\eps$, peak KL coef
$\beta_0$, mini-epochs $E$, filter bounds
$[\alpha_{\min}, \alpha_{\max}]$, Huber threshold $\delta$, lr schedule
$\eta_t$ with peak $\eta_{\max}$.
\begin{enumerate}[topsep=2pt,itemsep=1pt,leftmargin=*]
\item For each prompt $x \sim \D$:
  \begin{enumerate}[topsep=1pt,itemsep=1pt,label=\alph*.]
  \item \emph{Rollout} (train mode, dropout on). For $k=1,\dots,K$:
    sample fresh dropout RNG seed for $\xi^{(k)}$, store seed; capture
    chain RNG state; unroll latent recurrence
    \eqref{eq:dropout-recurrence} restoring the chain seed at every
    latent step (locks $m_t \equiv m$); decode $y^{(k)}$;
    record per-token logprobs
    $\log\pi_{\theta_\old}(y^{(k)}\mid x,\xi^{(k)})$;
    compute $r^{(k)} = \mathrm{verifier}(y^{(k)}, x)$.
  \item \emph{Group filter (V7).} Compute the empirical group accuracy
        $\mu_r = \tfrac1K\sum_k r^{(k)}$.  If
        $\alpha_{\min} \le \mu_r \le \alpha_{\max}$, keep the group;
        otherwise drop it from this step's update set.  Coordinate the
        decision across DDP ranks via \texttt{all\_reduce}: if no group
        passes anywhere, skip the entire optimizer / scheduler step
        for this batch and proceed to the next.
  \item Compute advantages
        $A^{(k)} = r^{(k)} - \mu_r$ on the kept groups
        (\emph{no $\sigma_r$ normalization}; cf.\ Remark~\ref{rem:sigma}).
  \item \emph{Update} (mask replay). For mini-epoch $e=1,\dots,E$ and
        each kept $k$: regenerate $\xi^{(k)}$ from stored seed
        (RNG-restore at every latent step), recompute
        $\log\pi_\theta(y^{(k)}\mid x, \xi^{(k)})$, form
        $\rho^{(k)} = \pi_\theta/\pi_{\theta_\old}$, take a gradient
        step on the loss
        \begin{equation*}
        \Loss(\theta) \;=\;
          -\,\tfrac{1}{K_{\mathrm{kept}}}\!\sum_{k\,\mathrm{kept}}\!
            \tfrac{1}{|y^{(k)}|}\!\sum_{\ell}
            \min\!\big(\rho_\ell^{(k)} A^{(k)},\,
              \mathrm{clip}(\rho_\ell^{(k)},1\pm\eps) A^{(k)}\big)
          \;+\; \beta(t)\,\tilde k_3(r^{\mathrm{ref}})
        \end{equation*}
        where the per-rollout token sum is divided by the rollout
        length (V6, sequence-mean), $\beta(t) = \beta_0\,\eta_t/\eta_{\max}$
        anneals with the lr schedule (V5), and $\tilde k_3$ is the
        Huberized $k_3$ KL estimator of~\eqref{eq:huber-kl} below (V4).
  \item Set $\theta_\old \gets \theta$.
  \end{enumerate}
\end{enumerate}
\end{minipage}}
\label{alg:dropout-grpo}
\end{figure}

% =====================================================================
\section{Implementation refinements}
\label{sec:impl}
% =====================================================================

The theory in Sections~\ref{sec:method}--\ref{sec:mc-dropout} is unchanged
by the four implementation details below; each was adopted to address a
specific empirical pathology observed during training.

\subsection{Huberized $k_3$ KL estimator (V4)}

The $k_3$ KL estimator (Schulman, blog) used in the surrogate is
$k_3(r) = \exp(r) - r - 1$ with $r := \log\pi_\theta - \log\pi_{\mathrm{ref}}$.
A standard implementation choice is to clamp $r$ to $[-\delta, \delta]$
before evaluating $k_3$, on the grounds that $\exp(r)$ overflows for large
$r$.  This has the unintended effect of zeroing the gradient outside
$[-\delta, \delta]$ (the autograd derivative of \texttt{clamp} is zero
beyond its bounds), so any token whose log-ratio drifts past $\pm\delta$
escapes KL regularization entirely.  Empirically this is the dominant
pre-collapse signature: a small tail of escaped tokens drifts further with
no restoring force, eventually dragging the whole policy with it.

We replace the clamp with a Huber-style extension: $k_3$ inside, linear
extrapolation outside, value- and slope-matched at $\pm\delta$:
\begin{equation}
\tilde k_3(r) \;=\;
\begin{cases}
\exp(r) - r - 1
  & |r| \le \delta, \\[4pt]
(\mathrm{e}^{\delta} - 1)(r - \delta) + (\mathrm{e}^{\delta} - \delta - 1)
  & r > \delta, \\[4pt]
(\mathrm{e}^{-\delta} - 1)(r + \delta) + (\mathrm{e}^{-\delta} + \delta - 1)
  & r < -\delta.
\end{cases}
\label{eq:huber-kl}
\end{equation}
By construction $\tilde k_3 \in C^1(\R)$, $\tilde k_3 \ge 0$, and the
gradient $\tilde k_3'(r)$ is bounded in magnitude by $\mathrm{e}^{\delta} - 1$
in the right tail (and by $|\mathrm{e}^{-\delta} - 1| < 1$ in the left
tail) but is \emph{nonzero} everywhere outside $r = 0$.  Drifted tokens
therefore continue to feel a restoring pull, while the maximum per-token
contribution stays bounded.  In practice we use $\delta = 5$, giving a
maximum slope of $\mathrm{e}^5 - 1 \approx 147$.

\subsection{Trust-region annealing of $\beta$ (V5)}

With the lr schedule cosine-annealing from $\eta_{\max}$ to $\eta_{\min}$
post-warmup, holding $\beta$ fixed produces a slowly worsening dynamic:
the per-step policy update shrinks with $\eta_t$, but the per-step KL
pull is $\eta_t \cdot \beta \cdot \nabla\KL$, also linear in $\eta_t$, so
the \emph{relative} KL pull stays constant.  However, the policy gradient
direction integrates informatively across steps while the KL gradient
adds coherently toward $\theta_{\mathrm{ref}}$; once the policy reaches
the equilibrium where these magnitudes balance, further updates are pinned.

We anneal $\beta$ proportionally with $\eta_t$:
\begin{equation}
\beta(t) \;=\; \beta_0 \cdot \frac{\eta_t}{\eta_{\max}},
\label{eq:beta-anneal}
\end{equation}
so the per-update KL pull scales with $\eta_t^2$ rather than $\eta_t$.
This shifts the trust-region equilibrium outward as the lr decays, letting
the policy continue to drift away from $\pi_{\mathrm{ref}}$ on the slow
schedule.  During warmup, $\eta_t$ rises from near zero to $\eta_{\max}$,
and $\beta(t)$ rises with it: the policy has near-zero KL pull early on,
allowing free movement before the regularizer engages.

\subsection{Sequence-mean loss aggregation (V6)}

The original Algorithm~1 averaged $\Loss^{\mathrm{GRPO}}$ over rollouts but
left the per-rollout token aggregation unspecified.  The natural reading
matches the per-rollout normalization in
Algorithm~\ref{alg:dropout-grpo}: the per-token loss within rollout $k$
is averaged over its $|y^{(k)}|$ response tokens, and the resulting
per-rollout scalars are then averaged across $k$.  An alternative
``token-mean'' aggregation (sum all per-token losses, divide by total
token count across the batch) implicitly weights each rollout by its
length, entangling per-rollout influence with response length and
underweighting short responses.

For binary verifier rewards on math benchmarks, response lengths can vary
several-fold (concise correct answers vs.\ long incorrect ones), so the
two aggregations give visibly different gradient directions.  We use
sequence-mean throughout.  At $\theta = \theta_{\mathrm{old}}$ the policy
term reduces to $-A^{(k)}$ per rollout, and after summing across kept
rollouts in a group its expectation is exactly zero (by
$\sum_k A^{(k)} = 0$); the scalar loss displayed in training logs is
therefore dominated by the $\beta\,\tilde k_3$ term and is small in
absolute terms.  This is not a sign of vanishing gradient: the
per-rollout score $-A^{(k)}\,\nabla\log\pi_\theta(y^{(k)}\mid x,\xi^{(k)})$
is signed and uncorrelated across rollouts, so the gradient does not
cancel even when its scalar evaluation does.  Monitoring should use
$\|\nabla\Loss\|$ (or wandb \texttt{train/grad\_norm}) rather than the
loss value itself.

\subsection{Group-level accuracy filter and online curriculum (V7)}

The within-group reward standard deviation $\sigma_r$ has long been used
as a diagnostic for ``dead'' groups (Yu et al.\ 2024; ``Dynamic Sampling''
in DAPO).  We extend the diagnostic to a \emph{filter}: at the end of the
rollout phase, drop groups whose empirical accuracy
$\mu_r = \tfrac1K\sum_k r^{(k)}$ falls outside a configured window
$[\alpha_{\min}, \alpha_{\max}]$, where the defaults are $0.05$ and $0.95$
(a factor-of-two looser bound than DAPO's $\sigma_r > 0$, which
corresponds to $\mu_r \in (0, 1)$ but with hard endpoints).

The motivation is twofold.  First, groups at $\mu_r \approx 0$ (all
rollouts wrong, model can't yet solve) and $\mu_r \approx 1$ (all
rollouts right, model has mastered) carry no comparative information:
$A^{(k)} = r^{(k)} - \mu_r \approx 0$ for every $k$, so the per-token
gradient is essentially noise.  Second, dropping these groups
\emph{before} replay saves the corresponding theta-old / reference /
policy backward compute --- in our setup roughly $25\%$ of step time per
filtered group.

The filter can be tightened further (e.g.\ $0.15 \le \mu_r \le 0.85$) once
the policy is past the ``barely solving'' regime; this concentrates
compute on the ``marginal'' middle band where signal-to-noise is highest.
As the policy improves, the marginal band itself shifts: an online
curriculum emerges with no explicit data ordering.

\paragraph{Compatibility with Theorem~\ref{thm:unbiased}.}
The filter is a function of the rollout outcomes
$\{r^{(k)}\}_{k=1}^K$, which depend on $\theta$ through $y^{(k)}$.  It is
therefore not exogenous, and the unbiasedness statement of
Theorem~\ref{thm:unbiased} no longer applies to $\nabla J(\theta)$ on the
original prompt distribution.  What it does apply to is
$\nabla J_{\mathrm{filt}}(\theta) :=
   \E_{x \sim \D}\big[ \E_{(\xi,y)\sim\Pi_\theta}[r(y,x)] \;\big|\;
   \mu_r(x) \in [\alpha_{\min},\alpha_{\max}] \big]$,
i.e.\ the expected reward conditional on the filter accepting.  The
implicit reweighting of the prompt distribution toward the marginal
band is the curriculum effect, by design.  In particular, with
$\alpha_{\min} = 0$ and $\alpha_{\max} = 1$ the filter is exactly the
zero-variance mask, and Theorem~\ref{thm:unbiased} recovers up to the
zero-variance subset (where $A^{(k)} \equiv 0$ contributes nothing
anyway).

\paragraph{Skip coordination under DDP.}
DDP requires every rank to run \texttt{loss.backward()} in lockstep --- a
rank that skips its backward pass causes the others to hang on gradient
sync.  We coordinate filter decisions globally via
\texttt{dist.all\_reduce} on the per-rank kept-group count: if the global
count is zero, every rank skips the policy update for this batch (no
optimizer step, no scheduler step --- a true no-op step, advancing only
the data iterator and step counter); if the global count is positive but
some rank has zero kept groups locally, that rank falls back to using
all its groups (whose pure-KL gradient is then diluted by $1/W$ in the
DDP gradient average over $W$ ranks).  In the symmetric multi-rank
regime the dilution is bounded; in the asymmetric regime where it would
not be, the all-zero global case dominates and a true skip is taken.
The constraint
$\texttt{rollout\_chunk\_size} = K$ ensures each chunk corresponds to
exactly one group, so dropping a group is a clean structural operation
that preserves the chunk-aligned RNG-restore replay.

% =====================================================================
\section{Practical concerns and open questions}
% =====================================================================

\paragraph{Sensitivity of $h_T$ to $\xi$.}
If dropout perturbations to $h_T$ do not change the argmax of the
answer-token distribution, all $K$ rollouts agree, advantages collapse to
zero, and no gradient is produced.  This is the dropout-rate analogue of
temperature in token-level RL.  A diagnostic to monitor is the
within-group standard deviation of $r^{(k)}$ (used in our implementation
for zero-variance group filtering); if $\sigma_r \to 0$, raise $p$ or fold
token-sampling Gumbel noise into $\xi$.  Note that $\sigma_r$ is used here
only as a \emph{diagnostic / filter}, not in the advantage --- consistent
with Theorem~\ref{thm:unbiased}.

\paragraph{Importance-ratio explosion.}
Under mask replay, $\rho_\theta(\xi,y)$ depends only on the parameter
update $\theta - \theta_\old$ (the mask-induced logit shifts are common
to numerator and denominator), so the standard PPO Lipschitz argument
applies directly.  Without mask replay, $\rho$ would also reflect mask
resampling and could be large because dropping different units changes
logits non-smoothly --- this is the variance source that mask replay
removes (cf.\ Proposition~\ref{prop:crn}).

\paragraph{Validity of clip under CRN.}
The clip operator in~\eqref{eq:grpo-loss} is non-smooth, so
Theorem~\ref{thm:unbiased} as stated holds only at $\theta = \theta_\old$
(initial step within each mini-epoch).  Standard PPO analysis (Schulman
et al., 2017, \S 5) extends the bound to small steps under a Lipschitz
condition on $\rho_\theta$, which under mask replay reduces to
Lipschitzness of $\pi_\theta$ in $\theta$ at fixed $\xi$ --- a mild
condition implied by Assumption~\ref{ass:smooth} together with bounded
weight matrices.  Empirical monitoring of the maximum log-ratio is
recommended.

\paragraph{KL anchor.}
Theorem~\ref{thm:unbiased} retains the KL term in the gradient
($-\beta\nabla\KL$ contribution).  If one wishes the surrogate-only
unbiasedness to be a statement about the full GRPO loss, set
$\pi_{\mathrm{ref}} \gets \pi_{\theta_\old}$ at the start of each update
(reference renewal), so that $\nabla\KL|_{\theta_\old} = 0$.  Otherwise
the KL gradient is a known regularizer term and Theorem~\ref{thm:unbiased}
still pins the policy-gradient component exactly.

\paragraph{Train/deploy mismatch.}
Inference typically disables dropout, so the deployed policy is
$\pi_\theta(y\mid x)$ with no mask, rather than the marginal
$\bar\pi_\theta$.  Two options: (a)~keep dropout on at deploy time and
average over masks, or (b)~accept a bias whose magnitude depends on the
network's per-layer Lipschitz constants; we do not give a closed-form
bound here and recommend empirical comparison.

\paragraph{Memory for stored masks.}
Storing $\xi^{(k)}$ for $K$ rollouts costs $\mathcal{O}(K \cdot
\mathrm{seed\_size})$ bits when implemented via per-chain RNG seeds, since
the mask is regenerated deterministically at update time.  This is the
recommended implementation and the one used in
\texttt{coconut\_variational.py}.

\paragraph{Constant $(1 - 1/K)$ scaling.}
Theorem~\ref{thm:unbiased} introduces a $(1 - 1/K)$ multiplicative factor
on the policy gradient relative to the true $\nabla J$.  For $K \ge 4$
this is $\ge 0.75$ and is absorbed into the learning rate without
practical consequence.  For very small $K$ (e.g.\ $K = 2$ giving factor
$0.5$) the effective learning rate halves; bear this in mind when
sweeping group sizes.

% =====================================================================
\section{Summary}
% =====================================================================

\begin{itemize}[topsep=2pt,itemsep=2pt]
\item Latent-reasoning models break GRPO because the latent phase is
      deterministic; variational dropout (one mask per rollout, shared
      across all latent steps) supplies the missing stochasticity without
      changing the model class.
\item Dropout noise is exogenous (independent of $\theta$); the score
      function reduces to the $\xi$-conditional log likelihood
      (Eq.~\ref{eq:score-reduces}).
\item The latent gradient $\partial h_T/\partial \theta$ is well defined
      at every depth $T$ (Lemma~\ref{lem:induction}), justifying
      backprop through the unrolled latent recurrence.
\item GRPO with $A^{(k)} = r^{(k)} - \mu_r$ (no $\sigma_r$ normalization)
      and replayed masks has surrogate gradient equal to
      $(1 - 1/K)\,\nabla J(\theta)$ at $\theta = \theta_\old$
      (Theorem~\ref{thm:unbiased}, Corollary~\ref{cor:marginal}).  The
      $\sigma_r$ normalization is dropped because it introduces bias
      (Remark~\ref{rem:sigma}; cf.\ Liu et al.\ 2024).
\item Replaying the rollout-time mask gives a CRN coupling that
      lower-bounds variance against a fresh-rollout estimator
      (Proposition~\ref{prop:crn}), with the bound becoming tight in the
      trust-region regime.
\item The shared-mask construction makes the marginal policy a true
      Bayesian model average $\bar\pi_\theta(y\mid x)
        = \E_{\theta'\sim q_\theta}[\pi_{\theta'}(y\mid x)]$
      (Section~\ref{sec:mc-dropout}, Eq.~\ref{eq:bma}).
\item Four implementation refinements (Section~\ref{sec:impl}) stabilize
      training without altering the theoretical claims:
      Huberized $k_3$ KL with bounded but nonzero tail gradient
      (Eq.~\ref{eq:huber-kl});
      lr-coupled annealing of $\beta$
      (Eq.~\ref{eq:beta-anneal}) to keep the trust region proportional to
      the step size;
      sequence-mean loss aggregation in line with Algorithm~\ref{alg:dropout-grpo};
      and a group-level accuracy filter
      $\alpha_{\min} \le \mu_r \le \alpha_{\max}$ with collective DDP
      skip coordination, acting as an online curriculum that targets the
      marginal-difficulty band where signal-to-noise is highest.
\item Practical caveats are: $h_T$-sensitivity, importance-ratio control,
      clip-region Lipschitz behavior, KL anchor handling, and
      train/deploy mismatch.
\end{itemize}

% =====================================================================
\section*{References (informal)}
% =====================================================================

\begin{itemize}[topsep=2pt,itemsep=1pt,leftmargin=*]
\item Hao, S.\ et al.\ (2024). \emph{Training large language models to reason
      in a continuous latent space.}  (Coconut.)
\item Shao, Z.\ et al.\ (2024).  \emph{DeepSeekMath: Pushing the limits of
      mathematical reasoning in open language models.} (GRPO.)
\item Liu, Z.\ et al.\ (2024).  \emph{Understanding R1-Zero-like training:
      A critical perspective.} (Dr.\ GRPO --- bias of $\sigma_r$
      normalization.)
\item Yu, Q.\ et al.\ (2024).  \emph{DAPO: An open-source LLM reinforcement
      learning system at scale.}  (Dynamic sampling / zero-variance group
      filtering.)
\item Schulman, J.\ et al.\ (2017).  \emph{Proximal policy optimization
      algorithms.}
\item Schulman, J.\ et al.\ (2015).  \emph{Gradient estimation using
      stochastic computation graphs.}
\item Gal, Y.\ \& Ghahramani, Z.\ (2016).  \emph{Dropout as a Bayesian
      approximation: Representing model uncertainty in deep learning.}
\item Glasserman, P.\ \& Yao, D.D.\ (1992).  \emph{Some guidelines and
      guarantees for common random numbers.}
\end{itemize}

\end{document}
